%============================================================
% HALAMAN ABSTRAK  (ringkasan dalam Bahasa Indonesia)
% Isi: latar belakang, tujuan KP, kegiatan yang dilakukan,
%      hasil, dan manfaat yang diperoleh.
%============================================================
\cleardoublepage
\phantomsection
\addcontentsline{toc}{chapter}{ABSTRAK}
\begin{center}
  {\fontsize{12}{15}\bfseries ABSTRAK\par}
\end{center}
\vspace{12pt}

Laporan ini membahas pelaksanaan Kerja Praktik yang dilakukan di PT. Sinergi Dimensi
Informatika (SINDIKA), dengan fokus pada pengembangan sistem monitoring infrastruktur
kelistrikan `Voltra' dan library backend Nedo.AspNet.Request.Validation. Penulis terlibat
langsung dalam implementasi logika bisnis pemantauan aset kelistrikan pada sistem Voltra,
mencakup kalkulasi Health Index, Critical Index, dan Risk Index sebagai dasar klasifikasi
status aset (Good, Warning, Critical). Selain itu, penulis juga mengembangkan fitur
manajemen tiket insiden yang terintegrasi dengan mekanisme auto-close berbasis event,
modul analitik finansial untuk perhitungan cost avoidance dan risk cost projection,
evaluasi keusangan aset (obsolescence) beserta prediksi Remaining Useful Life (RUL),
serta sistem pelaporan otomatis berbasis PDF. Sistem Voltra dibangun menggunakan pendekatan arsitektur microservices dengan RabbitMQ
sebagai message broker untuk komunikasi asinkron, MQTT untuk ingesti data sensor dari
perangkat lapangan, Redis sebagai cache dan pub/sub engine, serta kombinasi InfluxDB dan
PostgreSQL untuk penyimpanan data time-series dan relasional. Di luar pengembangan Voltra,
penulis turut merancang dan mengembangkan library Nedo.AspNet.Request.Validation, sebuah
pustaka validasi input berbasis attribute untuk aplikasi ASP.NET Core yang bertujuan
menyediakan mekanisme validasi terpusat, konsisten, dan dapat digunakan kembali di seluruh
ekosistem backend SINDIKA, lengkap dengan standarisasi kode error dan dukungan multi-bahasa
(localization) pada respons validasi API. Melalui pendekatan profesional dan kolaboratif selama kerja praktik, kegiatan ini memberikan
pemahaman mendalam mengenai penerapan arsitektur microservices, komunikasi asinkron, serta
praktik pengembangan library internal dalam skenario industri nyata.

\vspace{0.8cm}
\noindent\textbf{Kata Kunci:} Voltra, Microservices, RabbitMQ, MQTT, Health Index, Risk Index, Library Validation, ASP.NET Core, Predictive Maintenance
